\section*{补充推导：固定事件下的热能耗表达}
\label{sec:formal-details}

本部分集中给出系数来源、条件约束的有效界与对应关系。推导使用精确实数系数，具体数值计算及误差协议留待实现阶段。

\subsection*{一、固定控制状态下的物理系数}

在联合状态\(\omega\)下，定义参考温度偏置与基准注热向量
\begin{equation}
(b^P_\omega)_n=
\sum_g a^B_{ng}(\xi_g(\omega))B_g(\xi_g(\omega))
+\sum_g\lambda_{gn}(\xi_g(\omega))(T_a-T_g^{\rm ref}).
\label{eq:baseline-vector}
\end{equation}
令\((Q_c)_{no}=h_{n,u_o}q_o\)，并令\(\ell_{\omega,n}=\sum_g\lambda_{gn}(\xi_g(\omega))\)。由各来源的分配守恒，有
\begin{equation}
\begin{aligned}
p^{\rm heat}&=b^P_\omega+\operatorname{diag}(\ell_\omega)\theta+Q_cv,\\
b_\omega&=C^{-1}b^P_\omega,\qquad D_c=C^{-1}Q_c,\\
\beta_\omega&=\mathbf1^{\mathsf T}b^P_\omega,\qquad
P^{\rm net}=\beta_\omega+\ell_\omega^{\mathsf T}\theta+q_c^{\mathsf T}v.
\end{aligned}
\label{eq:physical-coefficients}
\end{equation}
其中\(\mathbf1^{\mathsf T}Q_c=q_c^{\mathsf T}\)。\(A_\omega\)单位为\(\mathrm{s}^{-1}\)，\(b_\omega,D_c\)对无量纲活动选择产生K/s的温度变化率，\(\beta_\omega,q_c\)单位为W，\(\ell_\omega\)单位为W/K。式\eqref{eq:physical-coefficients}同时连接电功率和注热，避免两套评价使用不一致的功率。

\subsection*{二、温度传播与完整能量积分}

对任意固定片长\(\delta\)，定义
\begin{equation}
\Phi_\omega(\delta)=e^{A_\omega\delta},\quad
\mathsf K_\omega(\delta)=\int_0^\delta e^{A_\omega t}\,dt,\quad
\mathsf L_\omega(\delta)=\int_0^\delta(\delta-t)e^{A_\omega t}\,dt.
\label{eq:matrix-integrals}
\end{equation}
由于片段内\(v\)固定，温度解为
\begin{equation}
\theta(t)=e^{A_\omega t}\theta(0)
+\mathsf K_\omega(t)(b_\omega+D_cv).
\label{eq:segment-solution}
\end{equation}
再对电功率积分，交换积分次序可得
\begin{equation}
\begin{aligned}
\theta(\delta)&=\Phi_\omega(\delta)\theta(0)
+\mathsf K_\omega(\delta)b_\omega+\mathsf K_\omega(\delta)D_cv,\\
\Delta E&=\ell_\omega^{\mathsf T}\mathsf K_\omega(\delta)\theta(0)
+\beta_\omega\delta+\ell_\omega^{\mathsf T}\mathsf L_\omega(\delta)b_\omega\\
&\quad+\bigl[\delta q_c^{\mathsf T}
+\ell_\omega^{\mathsf T}\mathsf L_\omega(\delta)D_c\bigr]v.
\end{aligned}
\label{eq:segment-integrals}
\end{equation}
因此，正文中的系数分别为
\begin{equation}
\begin{aligned}
\psi&=\mathsf K b_\omega,& \Gamma&=\mathsf K D_c,&
\mu^{\mathsf T}&=\ell_\omega^{\mathsf T}\mathsf K,\\
\zeta&=\beta_\omega\delta+\ell_\omega^{\mathsf T}\mathsf Lb_\omega,&
\nu^{\mathsf T}&=\delta q_c^{\mathsf T}+\ell_\omega^{\mathsf T}\mathsf LD_c.
\end{aligned}
\label{eq:integral-coefficients}
\end{equation}
各\(\mathsf K,\mathsf L\)均在当前\(\omega,\delta_{c\ell}\)取值。该定义不使用\(A_\omega^{-1}\)，所以不要求矩阵可逆；有限时间上的传播也不要求其渐近稳定。积分中的活动功率、等待及温度反馈已经完整计入\(\Delta E\)。

\subsection*{三、有依据的变量界与条件线性化}

固定事件和有限状态集合使得所有系数有界。取相容的无穷范数，定义
\begin{equation}
a_* =\max_{\omega\in\Omega}\|A_\omega\|_\infty,\qquad
b_{c*}=\max_{\omega\in\Omega}
\bigl(\|b_\omega\|_\infty+\|D_c\|_\infty\bigr).
\label{eq:norm-coefficients}
\end{equation}
由于\(\|v\|_\infty\le1\)，分段轨迹满足
\(\|\dot\theta\|_\infty\le a_*\|\theta\|_\infty+b_{c*}\)。由积分不等式可得对整个窗口有效的保守界
\begin{equation}
\|\theta_c(t)\|_\infty\le H_c,
\qquad H_c=e^{a_*W_c}
\bigl(\|T_c^0-T_a\mathbf1\|_\infty+b_{c*}W_c\bigr).
\label{eq:state-bound}
\end{equation}
该界覆盖任意合法状态切换及选项活动，包括尚未满足温限的轨迹。它证明存在有限界，不代表该界适合直接用于高效数值求解。进一步令
\begin{equation}
\overline P_c=\max_\omega
\left(|\beta_\omega|+\|\ell_\omega\|_1H_c+\|q_c\|_1\right),
\qquad |\Delta E_{c\ell}|\le\delta_{c\ell}\overline P_c.
\label{eq:energy-bound}
\end{equation}
可在主问题中加入\(-H_c\le\theta_{c\ell n}\le H_c\)及上述段能量界。完整物理功率的非负性仍由有效域条件约束，绝对值界本身不改变计量口径。

对式\eqref{eq:thermal-transition}第\(n\)行及式\eqref{eq:energy-transition}，一组可用的残差界为
\begin{equation}
\begin{aligned}
M^T_n&=H_c+\|\Phi_{n:}\|_1H_c+|\psi_n|+\|\Gamma_{n:}\|_1,\\
M^E&=\delta_{c\ell}\overline P_c+\|\mu\|_1H_c+|\zeta|+\|\nu\|_1.
\end{aligned}
\label{eq:residual-bounds}
\end{equation}
温度残差使用式\eqref{eq:thermal-bigM}，能量残差采用相同的双边形式。\(w=1\)时等式生效，\(w=0\)时上述界覆盖全变量盒上的残差。实际实现可以根据必要条件收紧界，但须保留原有限域全部可行安排，不能以任意小常数替代。

有效温区以及组功率、节点注热非负等条件，在固定\(\omega\)下可统一写成
\begin{equation}
a_{j\omega}^{\mathsf T}\theta+b_{jc\omega}^{\mathsf T}v+d_{j\omega}\le0.
\label{eq:validity-affine}
\end{equation}
对每个片段，以当前片段的\(\omega,v\)分别约束其两端温度，并用\(w\)条件激活；这同时覆盖内部切换处左右两侧的功率状态。窗口末点\(W_c\)按半开执行区间已无活动，还须在最后保持状态下补充式\eqref{eq:validity-affine}的\((\theta_{cL_c},v=0)\)条件。该末点检查不改变能量积分。段内有效性仍需连续检查。若在\(\tau^*\)发现明确违反，将式\eqref{eq:interior-temperature}代入式\eqref{eq:validity-affine}即可得到必要反馈行。

一般的单边条件行\(w=1\Rightarrow a^{\mathsf T}\theta+b^{\mathsf T}v+d\le0\)，可取\(M=\|a\|_1H_c+\|b\|_1+|d|\)，写成\(a^{\mathsf T}\theta+b^{\mathsf T}v+d\le M(1-w)\)。这也给出了式\eqref{eq:conditional-cut}的有限线性化界。若未来采用不同的有效域，应重新核查其表达，不能由本式自动覆盖。

\subsection*{四、对应关系与连续检查的作用}

给定\(P_{\mathcal G}\)的一份合法安排，每个实例对应唯一的模式--起点选项，因而能构造满足任务分配、期限和占用约束的\(x,y\)。实际域状态给出\(\rho,w\)：持续Stage锁点，无新启动时保持，入口已知，因此控制约束成立。按真实热轨迹填写事件温度和段能量，积分等式成立，所有安全必要行也成立。

反过来，整数选项约束为每个实例选择一项，并将其单元绑定到共享部署。固定事件覆盖全部执行区间，资源互斥与控制约束恢复合法时间表及唯一保持状态；在给定入口下，线性微分方程有唯一解，条件传播与积分等式正是该解的端点和能量。若进一步通过完整连续检查，便得到\(P_{\mathcal G}\)的可行安排。由此，有限热行主问题与完整受限问题的差异在于尚未加入的连续必要条件，而不是模式或等待被提前删去。

对任一明确超温的候选，式\eqref{eq:conditional-cut}在其状态下被违反；对任一安全安排，该状态激活时必须满足该行，未激活时不受其限制。故该行保留所有安全安排。完整赋值排除也只删除已有违规证明的完整安排。前提是系数及违规判断可靠；未分辨不能充当排除证明。

一般多节点的温度导数是矩阵指数表达，其内部驻点不能直接沿用旧二节点闭式程序。实现时须提供覆盖全区间的根区间检查或温度包络，处理极值并列、切换端点及靠近温限的未分辨情形。优化系数的近似误差还需单独核查：最终连续检查可以检验一个候选，却不能恢复此前被错误系数排除的候选。数学最优结论与实际求解器报告状态因而分别记录。
